% paper17-fleet-competition.tex
% Fleet Competition: Adversarial Evidence Selection via Tournament Protocols
% Paper 17 in the Judgment Geometry series.
\documentclass[11pt]{article}
% jugeo-common.tex — shared preamble for all Judgment Geometry papers
% Include via: % jugeo-common.tex — shared preamble for all Judgment Geometry papers
% Include via: % jugeo-common.tex — shared preamble for all Judgment Geometry papers
% Include via: \input{jugeo-common}

% ── Packages ────────────────────────────────────────────────────────────
\usepackage{amsmath,amssymb,amsthm,mathtools}
\usepackage{tikz-cd}
\usepackage{listings}
\usepackage{xcolor}
\usepackage{xspace}
\usepackage{booktabs}
\usepackage{multirow}
\usepackage{enumitem}
\usepackage[expansion=false]{microtype}
\usepackage{hyperref}
\usepackage{cleveref}
% \usepackage{thmtools}  % disabled: not available in this TeX installation
\usepackage{float}
\usepackage{caption}
\usepackage{subcaption}
\usepackage{tabularx}
\usepackage{stmaryrd}       % \llbracket, \rrbracket
\usepackage{bbm}            % \mathbbm{1}

% ── Page geometry ───────────────────────────────────────────────────────
\usepackage[margin=1in]{geometry}

% ── Theorem environments ────────────────────────────────────────────────
\theoremstyle{plain}
\newtheorem{theorem}{Theorem}[section]
\newtheorem{lemma}[theorem]{Lemma}
\newtheorem{proposition}[theorem]{Proposition}
\newtheorem{corollary}[theorem]{Corollary}

\theoremstyle{definition}
\newtheorem{definition}[theorem]{Definition}
\newtheorem{example}[theorem]{Example}
\newtheorem{construction}[theorem]{Construction}

\theoremstyle{remark}
\newtheorem{remark}[theorem]{Remark}
\newtheorem{notation}[theorem]{Notation}

% ── Python listings style ──────────────────────────────────────────────
\definecolor{jugeoBlue}{HTML}{2563EB}
\definecolor{jugeoGreen}{HTML}{059669}
\definecolor{jugeoRed}{HTML}{DC2626}
\definecolor{jugeoPurple}{HTML}{7C3AED}
\definecolor{jugeoGray}{HTML}{6B7280}
\definecolor{jugeoBg}{HTML}{F8FAFC}

\lstdefinestyle{jugeo-python}{
  language=Python,
  basicstyle=\ttfamily\small,
  keywordstyle=\color{jugeoBlue}\bfseries,
  stringstyle=\color{jugeoGreen},
  commentstyle=\color{jugeoGray}\itshape,
  numberstyle=\tiny\color{jugeoGray},
  backgroundcolor=\color{jugeoBg},
  numbers=left,
  numbersep=6pt,
  frame=single,
  framerule=0.4pt,
  rulecolor=\color{jugeoGray!40},
  breaklines=true,
  breakatwhitespace=true,
  showstringspaces=false,
  tabsize=4,
  xleftmargin=1.5em,
  framexleftmargin=1.5em,
  morekeywords={dataclass,frozen,field,Enum,IntEnum,auto,Any,
                Mapping,Sequence,Optional,match,case},
  literate={→}{{$\to$}}1 {←}{{$\leftarrow$}}1
           {≤}{{$\leq$}}1 {≥}{{$\geq$}}1 {≠}{{$\neq$}}1
           {∈}{{$\in$}}1 {∀}{{$\forall$}}1 {∃}{{$\exists$}}1
           {⊕}{{$\oplus$}}1 {⊖}{{$\ominus$}}1,
  escapeinside={(*@}{@*)},
}
\lstset{style=jugeo-python}

\lstdefinestyle{jugeo-output}{
  basicstyle=\ttfamily\small,
  backgroundcolor=\color{jugeoBg},
  frame=single,
  framerule=0.4pt,
  rulecolor=\color{jugeoGray!40},
  breaklines=true,
  showstringspaces=false,
  xleftmargin=1.5em,
  framexleftmargin=0.5em,
  numbers=none,
}

% ── JuGeo-specific macros ──────────────────────────────────────────────

% Judgment 8-tuple
\newcommand{\J}{\mathcal{J}}
\newcommand{\Jud}[8]{(#1,\,#2,\,#3,\,#4,\,#5,\,#6,\,#7,\,#8)}
\newcommand{\JudShort}[1]{\J_{#1}}

% Components
\newcommand{\coord}{\mathsf{c}}            % coordinate
\newcommand{\prop}{\varphi}                 % proposition
\newcommand{\carrier}{\mathcal{A}}          % carrier type
\newcommand{\evid}{\mathcal{E}}             % evidence bundle
\newcommand{\oblig}{\mathcal{O}}            % obligations
\newcommand{\obstr}{\mathcal{B}}            % obstructions
\newcommand{\trust}{\mathcal{T}}            % trust annotation
\newcommand{\prov}{\Pi}                     % provenance

% Trust algebra
\newcommand{\Talg}{\mathfrak{T}}            % trust ordered algebra
\newcommand{\Eadm}{\mathcal{E}_{\mathrm{adm}}}
\newcommand{\tleq}{\preceq}                 % trust ordering
\newcommand{\tjoin}{\oplus}                 % conservative join
\newcommand{\tatten}{\ominus}               % attenuation
\newcommand{\tpromote}{\uparrow_\pi}        % promotion
\newcommand{\tdemote}{\downarrow_\chi}       % demotion

% Trust tiers
\newcommand{\Tcontra}{\ensuremath{\mathsf{CONTRADICTED}}}
\newcommand{\Tunver}{\ensuremath{\mathsf{UNVERIFIED}}}
\newcommand{\Tcopilot}{\ensuremath{\mathsf{COPILOT}}}
\newcommand{\Toracle}{\ensuremath{\mathsf{ORACLE}}}
\newcommand{\Truntime}{\ensuremath{\mathsf{RUNTIME}}}
\newcommand{\Tsolver}{\ensuremath{\mathsf{SOLVER}}}
\newcommand{\Tproof}{\ensuremath{\mathsf{PROOF}}}

% Site and geometry
\newcommand{\Site}{\mathcal{S}}             % semantic site
\newcommand{\Cov}{\mathrm{Cov}}            % covering family
\newcommand{\Ob}{\mathrm{Ob}}              % objects
\newcommand{\Mor}{\mathrm{Mor}}            % morphisms
\newcommand{\res}{\mathrm{res}}            % restriction
\newcommand{\Cech}{\check{C}}              % Čech complex
\newcommand{\Hcoh}[1]{H^{#1}}             % cohomology group

% Descent
\newcommand{\descent}{\mathrm{desc}}
\newcommand{\glue}{\mathrm{glue}}
\newcommand{\overlap}[2]{#1 \cap #2}

% Semantic moves
\newcommand{\VERIFY}{\textsc{Verify}}
\newcommand{\CONSTRUCT}{\textsc{Construct}}
\newcommand{\NORMALIZE}{\textsc{Normalize}}
\newcommand{\GLUE}{\textsc{Glue}}
\newcommand{\RESTRICT}{\textsc{Restrict}}
\newcommand{\TRANSPORT}{\textsc{Transport}}
\newcommand{\REPAIR}{\textsc{Repair}}
\newcommand{\PROMOTE}{\textsc{Promote}}
\newcommand{\CHALLENGE}{\textsc{Challenge}}
\newcommand{\ESCALATE}{\textsc{Escalate}}

% Fragments
\newcommand{\QFLIA}{\texttt{QF\_LIA}}
\newcommand{\QFLRA}{\texttt{QF\_LRA}}
\newcommand{\QFBV}{\texttt{QF\_BV}}
\newcommand{\QFUF}{\texttt{QF\_UF}}

% Misc
\newcommand{\Python}{\textsc{Python}}
\newcommand{\JuGeo}{\textsc{JuGeo}}
\newcommand{\LEAN}{\textsc{Lean}\,4}
\newcommand{\Fstar}{F$^{\star}$}
\newcommand{\Coq}{\textsc{Coq}}
\newcommand{\Dafny}{\textsc{Dafny}}

% Common shortcuts
\newcommand{\ie}{i.e.\@\xspace}
\newcommand{\eg}{e.g.\@\xspace}
\newcommand{\cf}{cf.\@\xspace}
\newcommand{\wrt}{w.r.t.\@\xspace}

% Evaluation shortcuts
\newcommand{\numcases}{300}
\newcommand{\numspec}{100}
\newcommand{\numeq}{100}
\newcommand{\numbug}{100}
\newcommand{\medtime}{6.5\,ms}
\newcommand{\totaltime}{1.949\,s}


% ── Packages ────────────────────────────────────────────────────────────
\usepackage{amsmath,amssymb,amsthm,mathtools}
\usepackage{tikz-cd}
\usepackage{listings}
\usepackage{xcolor}
\usepackage{xspace}
\usepackage{booktabs}
\usepackage{multirow}
\usepackage{enumitem}
\usepackage[expansion=false]{microtype}
\usepackage{hyperref}
\usepackage{cleveref}
% \usepackage{thmtools}  % disabled: not available in this TeX installation
\usepackage{float}
\usepackage{caption}
\usepackage{subcaption}
\usepackage{tabularx}
\usepackage{stmaryrd}       % \llbracket, \rrbracket
\usepackage{bbm}            % \mathbbm{1}

% ── Page geometry ───────────────────────────────────────────────────────
\usepackage[margin=1in]{geometry}

% ── Theorem environments ────────────────────────────────────────────────
\theoremstyle{plain}
\newtheorem{theorem}{Theorem}[section]
\newtheorem{lemma}[theorem]{Lemma}
\newtheorem{proposition}[theorem]{Proposition}
\newtheorem{corollary}[theorem]{Corollary}

\theoremstyle{definition}
\newtheorem{definition}[theorem]{Definition}
\newtheorem{example}[theorem]{Example}
\newtheorem{construction}[theorem]{Construction}

\theoremstyle{remark}
\newtheorem{remark}[theorem]{Remark}
\newtheorem{notation}[theorem]{Notation}

% ── Python listings style ──────────────────────────────────────────────
\definecolor{jugeoBlue}{HTML}{2563EB}
\definecolor{jugeoGreen}{HTML}{059669}
\definecolor{jugeoRed}{HTML}{DC2626}
\definecolor{jugeoPurple}{HTML}{7C3AED}
\definecolor{jugeoGray}{HTML}{6B7280}
\definecolor{jugeoBg}{HTML}{F8FAFC}

\lstdefinestyle{jugeo-python}{
  language=Python,
  basicstyle=\ttfamily\small,
  keywordstyle=\color{jugeoBlue}\bfseries,
  stringstyle=\color{jugeoGreen},
  commentstyle=\color{jugeoGray}\itshape,
  numberstyle=\tiny\color{jugeoGray},
  backgroundcolor=\color{jugeoBg},
  numbers=left,
  numbersep=6pt,
  frame=single,
  framerule=0.4pt,
  rulecolor=\color{jugeoGray!40},
  breaklines=true,
  breakatwhitespace=true,
  showstringspaces=false,
  tabsize=4,
  xleftmargin=1.5em,
  framexleftmargin=1.5em,
  morekeywords={dataclass,frozen,field,Enum,IntEnum,auto,Any,
                Mapping,Sequence,Optional,match,case},
  literate={→}{{$\to$}}1 {←}{{$\leftarrow$}}1
           {≤}{{$\leq$}}1 {≥}{{$\geq$}}1 {≠}{{$\neq$}}1
           {∈}{{$\in$}}1 {∀}{{$\forall$}}1 {∃}{{$\exists$}}1
           {⊕}{{$\oplus$}}1 {⊖}{{$\ominus$}}1,
  escapeinside={(*@}{@*)},
}
\lstset{style=jugeo-python}

\lstdefinestyle{jugeo-output}{
  basicstyle=\ttfamily\small,
  backgroundcolor=\color{jugeoBg},
  frame=single,
  framerule=0.4pt,
  rulecolor=\color{jugeoGray!40},
  breaklines=true,
  showstringspaces=false,
  xleftmargin=1.5em,
  framexleftmargin=0.5em,
  numbers=none,
}

% ── JuGeo-specific macros ──────────────────────────────────────────────

% Judgment 8-tuple
\newcommand{\J}{\mathcal{J}}
\newcommand{\Jud}[8]{(#1,\,#2,\,#3,\,#4,\,#5,\,#6,\,#7,\,#8)}
\newcommand{\JudShort}[1]{\J_{#1}}

% Components
\newcommand{\coord}{\mathsf{c}}            % coordinate
\newcommand{\prop}{\varphi}                 % proposition
\newcommand{\carrier}{\mathcal{A}}          % carrier type
\newcommand{\evid}{\mathcal{E}}             % evidence bundle
\newcommand{\oblig}{\mathcal{O}}            % obligations
\newcommand{\obstr}{\mathcal{B}}            % obstructions
\newcommand{\trust}{\mathcal{T}}            % trust annotation
\newcommand{\prov}{\Pi}                     % provenance

% Trust algebra
\newcommand{\Talg}{\mathfrak{T}}            % trust ordered algebra
\newcommand{\Eadm}{\mathcal{E}_{\mathrm{adm}}}
\newcommand{\tleq}{\preceq}                 % trust ordering
\newcommand{\tjoin}{\oplus}                 % conservative join
\newcommand{\tatten}{\ominus}               % attenuation
\newcommand{\tpromote}{\uparrow_\pi}        % promotion
\newcommand{\tdemote}{\downarrow_\chi}       % demotion

% Trust tiers
\newcommand{\Tcontra}{\ensuremath{\mathsf{CONTRADICTED}}}
\newcommand{\Tunver}{\ensuremath{\mathsf{UNVERIFIED}}}
\newcommand{\Tcopilot}{\ensuremath{\mathsf{COPILOT}}}
\newcommand{\Toracle}{\ensuremath{\mathsf{ORACLE}}}
\newcommand{\Truntime}{\ensuremath{\mathsf{RUNTIME}}}
\newcommand{\Tsolver}{\ensuremath{\mathsf{SOLVER}}}
\newcommand{\Tproof}{\ensuremath{\mathsf{PROOF}}}

% Site and geometry
\newcommand{\Site}{\mathcal{S}}             % semantic site
\newcommand{\Cov}{\mathrm{Cov}}            % covering family
\newcommand{\Ob}{\mathrm{Ob}}              % objects
\newcommand{\Mor}{\mathrm{Mor}}            % morphisms
\newcommand{\res}{\mathrm{res}}            % restriction
\newcommand{\Cech}{\check{C}}              % Čech complex
\newcommand{\Hcoh}[1]{H^{#1}}             % cohomology group

% Descent
\newcommand{\descent}{\mathrm{desc}}
\newcommand{\glue}{\mathrm{glue}}
\newcommand{\overlap}[2]{#1 \cap #2}

% Semantic moves
\newcommand{\VERIFY}{\textsc{Verify}}
\newcommand{\CONSTRUCT}{\textsc{Construct}}
\newcommand{\NORMALIZE}{\textsc{Normalize}}
\newcommand{\GLUE}{\textsc{Glue}}
\newcommand{\RESTRICT}{\textsc{Restrict}}
\newcommand{\TRANSPORT}{\textsc{Transport}}
\newcommand{\REPAIR}{\textsc{Repair}}
\newcommand{\PROMOTE}{\textsc{Promote}}
\newcommand{\CHALLENGE}{\textsc{Challenge}}
\newcommand{\ESCALATE}{\textsc{Escalate}}

% Fragments
\newcommand{\QFLIA}{\texttt{QF\_LIA}}
\newcommand{\QFLRA}{\texttt{QF\_LRA}}
\newcommand{\QFBV}{\texttt{QF\_BV}}
\newcommand{\QFUF}{\texttt{QF\_UF}}

% Misc
\newcommand{\Python}{\textsc{Python}}
\newcommand{\JuGeo}{\textsc{JuGeo}}
\newcommand{\LEAN}{\textsc{Lean}\,4}
\newcommand{\Fstar}{F$^{\star}$}
\newcommand{\Coq}{\textsc{Coq}}
\newcommand{\Dafny}{\textsc{Dafny}}

% Common shortcuts
\newcommand{\ie}{i.e.\@\xspace}
\newcommand{\eg}{e.g.\@\xspace}
\newcommand{\cf}{cf.\@\xspace}
\newcommand{\wrt}{w.r.t.\@\xspace}

% Evaluation shortcuts
\newcommand{\numcases}{300}
\newcommand{\numspec}{100}
\newcommand{\numeq}{100}
\newcommand{\numbug}{100}
\newcommand{\medtime}{6.5\,ms}
\newcommand{\totaltime}{1.949\,s}


% ── Packages ────────────────────────────────────────────────────────────
\usepackage{amsmath,amssymb,amsthm,mathtools}
\usepackage{tikz-cd}
\usepackage{listings}
\usepackage{xcolor}
\usepackage{xspace}
\usepackage{booktabs}
\usepackage{multirow}
\usepackage{enumitem}
\usepackage[expansion=false]{microtype}
\usepackage{hyperref}
\usepackage{cleveref}
% \usepackage{thmtools}  % disabled: not available in this TeX installation
\usepackage{float}
\usepackage{caption}
\usepackage{subcaption}
\usepackage{tabularx}
\usepackage{stmaryrd}       % \llbracket, \rrbracket
\usepackage{bbm}            % \mathbbm{1}

% ── Page geometry ───────────────────────────────────────────────────────
\usepackage[margin=1in]{geometry}

% ── Theorem environments ────────────────────────────────────────────────
\theoremstyle{plain}
\newtheorem{theorem}{Theorem}[section]
\newtheorem{lemma}[theorem]{Lemma}
\newtheorem{proposition}[theorem]{Proposition}
\newtheorem{corollary}[theorem]{Corollary}

\theoremstyle{definition}
\newtheorem{definition}[theorem]{Definition}
\newtheorem{example}[theorem]{Example}
\newtheorem{construction}[theorem]{Construction}

\theoremstyle{remark}
\newtheorem{remark}[theorem]{Remark}
\newtheorem{notation}[theorem]{Notation}

% ── Python listings style ──────────────────────────────────────────────
\definecolor{jugeoBlue}{HTML}{2563EB}
\definecolor{jugeoGreen}{HTML}{059669}
\definecolor{jugeoRed}{HTML}{DC2626}
\definecolor{jugeoPurple}{HTML}{7C3AED}
\definecolor{jugeoGray}{HTML}{6B7280}
\definecolor{jugeoBg}{HTML}{F8FAFC}

\lstdefinestyle{jugeo-python}{
  language=Python,
  basicstyle=\ttfamily\small,
  keywordstyle=\color{jugeoBlue}\bfseries,
  stringstyle=\color{jugeoGreen},
  commentstyle=\color{jugeoGray}\itshape,
  numberstyle=\tiny\color{jugeoGray},
  backgroundcolor=\color{jugeoBg},
  numbers=left,
  numbersep=6pt,
  frame=single,
  framerule=0.4pt,
  rulecolor=\color{jugeoGray!40},
  breaklines=true,
  breakatwhitespace=true,
  showstringspaces=false,
  tabsize=4,
  xleftmargin=1.5em,
  framexleftmargin=1.5em,
  morekeywords={dataclass,frozen,field,Enum,IntEnum,auto,Any,
                Mapping,Sequence,Optional,match,case},
  literate={→}{{$\to$}}1 {←}{{$\leftarrow$}}1
           {≤}{{$\leq$}}1 {≥}{{$\geq$}}1 {≠}{{$\neq$}}1
           {∈}{{$\in$}}1 {∀}{{$\forall$}}1 {∃}{{$\exists$}}1
           {⊕}{{$\oplus$}}1 {⊖}{{$\ominus$}}1,
  escapeinside={(*@}{@*)},
}
\lstset{style=jugeo-python}

\lstdefinestyle{jugeo-output}{
  basicstyle=\ttfamily\small,
  backgroundcolor=\color{jugeoBg},
  frame=single,
  framerule=0.4pt,
  rulecolor=\color{jugeoGray!40},
  breaklines=true,
  showstringspaces=false,
  xleftmargin=1.5em,
  framexleftmargin=0.5em,
  numbers=none,
}

% ── JuGeo-specific macros ──────────────────────────────────────────────

% Judgment 8-tuple
\newcommand{\J}{\mathcal{J}}
\newcommand{\Jud}[8]{(#1,\,#2,\,#3,\,#4,\,#5,\,#6,\,#7,\,#8)}
\newcommand{\JudShort}[1]{\J_{#1}}

% Components
\newcommand{\coord}{\mathsf{c}}            % coordinate
\newcommand{\prop}{\varphi}                 % proposition
\newcommand{\carrier}{\mathcal{A}}          % carrier type
\newcommand{\evid}{\mathcal{E}}             % evidence bundle
\newcommand{\oblig}{\mathcal{O}}            % obligations
\newcommand{\obstr}{\mathcal{B}}            % obstructions
\newcommand{\trust}{\mathcal{T}}            % trust annotation
\newcommand{\prov}{\Pi}                     % provenance

% Trust algebra
\newcommand{\Talg}{\mathfrak{T}}            % trust ordered algebra
\newcommand{\Eadm}{\mathcal{E}_{\mathrm{adm}}}
\newcommand{\tleq}{\preceq}                 % trust ordering
\newcommand{\tjoin}{\oplus}                 % conservative join
\newcommand{\tatten}{\ominus}               % attenuation
\newcommand{\tpromote}{\uparrow_\pi}        % promotion
\newcommand{\tdemote}{\downarrow_\chi}       % demotion

% Trust tiers
\newcommand{\Tcontra}{\ensuremath{\mathsf{CONTRADICTED}}}
\newcommand{\Tunver}{\ensuremath{\mathsf{UNVERIFIED}}}
\newcommand{\Tcopilot}{\ensuremath{\mathsf{COPILOT}}}
\newcommand{\Toracle}{\ensuremath{\mathsf{ORACLE}}}
\newcommand{\Truntime}{\ensuremath{\mathsf{RUNTIME}}}
\newcommand{\Tsolver}{\ensuremath{\mathsf{SOLVER}}}
\newcommand{\Tproof}{\ensuremath{\mathsf{PROOF}}}

% Site and geometry
\newcommand{\Site}{\mathcal{S}}             % semantic site
\newcommand{\Cov}{\mathrm{Cov}}            % covering family
\newcommand{\Ob}{\mathrm{Ob}}              % objects
\newcommand{\Mor}{\mathrm{Mor}}            % morphisms
\newcommand{\res}{\mathrm{res}}            % restriction
\newcommand{\Cech}{\check{C}}              % Čech complex
\newcommand{\Hcoh}[1]{H^{#1}}             % cohomology group

% Descent
\newcommand{\descent}{\mathrm{desc}}
\newcommand{\glue}{\mathrm{glue}}
\newcommand{\overlap}[2]{#1 \cap #2}

% Semantic moves
\newcommand{\VERIFY}{\textsc{Verify}}
\newcommand{\CONSTRUCT}{\textsc{Construct}}
\newcommand{\NORMALIZE}{\textsc{Normalize}}
\newcommand{\GLUE}{\textsc{Glue}}
\newcommand{\RESTRICT}{\textsc{Restrict}}
\newcommand{\TRANSPORT}{\textsc{Transport}}
\newcommand{\REPAIR}{\textsc{Repair}}
\newcommand{\PROMOTE}{\textsc{Promote}}
\newcommand{\CHALLENGE}{\textsc{Challenge}}
\newcommand{\ESCALATE}{\textsc{Escalate}}

% Fragments
\newcommand{\QFLIA}{\texttt{QF\_LIA}}
\newcommand{\QFLRA}{\texttt{QF\_LRA}}
\newcommand{\QFBV}{\texttt{QF\_BV}}
\newcommand{\QFUF}{\texttt{QF\_UF}}

% Misc
\newcommand{\Python}{\textsc{Python}}
\newcommand{\JuGeo}{\textsc{JuGeo}}
\newcommand{\LEAN}{\textsc{Lean}\,4}
\newcommand{\Fstar}{F$^{\star}$}
\newcommand{\Coq}{\textsc{Coq}}
\newcommand{\Dafny}{\textsc{Dafny}}

% Common shortcuts
\newcommand{\ie}{i.e.\@\xspace}
\newcommand{\eg}{e.g.\@\xspace}
\newcommand{\cf}{cf.\@\xspace}
\newcommand{\wrt}{w.r.t.\@\xspace}

% Evaluation shortcuts
\newcommand{\numcases}{300}
\newcommand{\numspec}{100}
\newcommand{\numeq}{100}
\newcommand{\numbug}{100}
\newcommand{\medtime}{6.5\,ms}
\newcommand{\totaltime}{1.949\,s}

% experiment-data.tex — AUTO-GENERATED from real jugeo CLI runs
% DO NOT EDIT — regenerate with: python3 experiments/run_universal_metrics.py
% Generated from 30 programs across 6 domains

\newcommand{\expTotalPrograms}{30}
\newcommand{\expVerified}{30}
\newcommand{\expAccuracy}{100.0\%}
\newcommand{\expCoordsMin}{2}
\newcommand{\expCoordsMax}{10}
\newcommand{\expCoordsMean}{5.2}
\newcommand{\expCoordsMedian}{5.5}
\newcommand{\expPropsMin}{4}
\newcommand{\expPropsMax}{20}
\newcommand{\expPropsMean}{10.4}
\newcommand{\expPropsSum}{311}
\newcommand{\expPropsOkSum}{310}
\newcommand{\expTimeMin}{3.506\,s}
\newcommand{\expTimeMax}{6.714\,s}
\newcommand{\expTimeMean}{4.64\,s}
\newcommand{\expTimeMedian}{4.508\,s}
\newcommand{\expTimeTotal}{139.204\,s}
\newcommand{\expObstructionTotal}{0}
\newcommand{\expHOneZero}{30}
\newcommand{\expDomainCount}{6}
\newcommand{\expTrustCOPILOTSUGGESTED}{30}
\newcommand{\expDomainDsCount}{5}
\newcommand{\expDomainDsVerified}{5}
\newcommand{\expDomainDsAvgCoords}{7.6}
\newcommand{\expDomainDsAvgProps}{15.2}
\newcommand{\expDomainMathCount}{5}
\newcommand{\expDomainMathVerified}{5}
\newcommand{\expDomainMathAvgCoords}{4.8}
\newcommand{\expDomainMathAvgProps}{9.4}
\newcommand{\expDomainSmCount}{5}
\newcommand{\expDomainSmVerified}{5}
\newcommand{\expDomainSmAvgCoords}{7.6}
\newcommand{\expDomainSmAvgProps}{15.2}
\newcommand{\expDomainSortCount}{5}
\newcommand{\expDomainSortVerified}{5}
\newcommand{\expDomainSortAvgCoords}{2.4}
\newcommand{\expDomainSortAvgProps}{4.8}
\newcommand{\expDomainStrCount}{5}
\newcommand{\expDomainStrVerified}{5}
\newcommand{\expDomainStrAvgCoords}{3.2}
\newcommand{\expDomainStrAvgProps}{6.4}
\newcommand{\expDomainWebCount}{5}
\newcommand{\expDomainWebVerified}{5}
\newcommand{\expDomainWebAvgCoords}{5.6}
\newcommand{\expDomainWebAvgProps}{11.2}

% subsystem-data.tex — AUTO-GENERATED from real jugeo subsystem experiments
% DO NOT EDIT — regenerate with: python3 experiments/run_subsystem_experiments.py

\newcommand{\subCliTotal}{10}
\newcommand{\subCliVerified}{9}
\newcommand{\subCliAccuracy}{90.0\%}
\newcommand{\subCliMeanTime}{4.132\,s}
\newcommand{\subCliMinTime}{3.472\,s}
\newcommand{\subCliMaxTime}{5.372\,s}
\newcommand{\subCliTotalCoords}{54}
\newcommand{\subCliTotalProps}{107}
\newcommand{\subCliTotalPropsOk}{106}
\newcommand{\subSiteCalls}{90}
\newcommand{\subSiteSuccessful}{69}
\newcommand{\subSiteSuccessRate}{76.7\%}
\newcommand{\subSiteMeanTime}{0.0001\,s}
\newcommand{\subSiteTotalTime}{0.005\,s}
\newcommand{\subSpecificationsatisfactionSmallTime}{0.0003\,s}
\newcommand{\subSpecificationsatisfactionMediumTime}{0.0001\,s}
\newcommand{\subSpecificationsatisfactionLargeTime}{0.0001\,s}
\newcommand{\subInhabitantfleetSmallTime}{0.0\,s}
\newcommand{\subInhabitantfleetMediumTime}{0.0\,s}
\newcommand{\subInhabitantfleetLargeTime}{0.0\,s}
\newcommand{\subTheoremecologySmallTime}{0.0\,s}
\newcommand{\subTheoremecologyMediumTime}{0.0\,s}
\newcommand{\subTheoremecologyLargeTime}{0.0\,s}
\newcommand{\subStatespaceexplorationSmallTime}{0.0\,s}
\newcommand{\subStatespaceexplorationMediumTime}{0.0\,s}
\newcommand{\subStatespaceexplorationLargeTime}{0.0\,s}
\newcommand{\subRepairsemanticsSmallTime}{0.0\,s}
\newcommand{\subRepairsemanticsMediumTime}{0.0\,s}
\newcommand{\subRepairsemanticsLargeTime}{0.0\,s}
\newcommand{\subReplaygluingSmallTime}{0.0\,s}
\newcommand{\subReplaygluingMediumTime}{0.0\,s}
\newcommand{\subReplaygluingLargeTime}{0.0\,s}
\newcommand{\subSemanticfuturesSmallTime}{0.0\,s}
\newcommand{\subSemanticfuturesMediumTime}{0.0\,s}
\newcommand{\subSemanticfuturesLargeTime}{0.0\,s}
\newcommand{\subHypercovertreatySmallTime}{0.0\,s}
\newcommand{\subHypercovertreatyMediumTime}{0.0\,s}
\newcommand{\subHypercovertreatyLargeTime}{0.0\,s}
\newcommand{\subEncodeforsolverSmallTime}{0.0026\,s}
\newcommand{\subEncodeforsolverMediumTime}{0.0002\,s}
\newcommand{\subEncodeforsolverLargeTime}{0.0001\,s}
\newcommand{\subInterfaceroutingSmallTime}{0.0\,s}
\newcommand{\subInterfaceroutingMediumTime}{0.0\,s}
\newcommand{\subInterfaceroutingLargeTime}{0.0\,s}
\newcommand{\subOrchestrateverificationSmallTime}{0.0002\,s}
\newcommand{\subOrchestrateverificationMediumTime}{0.0\,s}
\newcommand{\subOrchestrateverificationLargeTime}{0.0\,s}
\newcommand{\subRunfulldescentSmallTime}{0.0\,s}
\newcommand{\subRunfulldescentMediumTime}{0.0\,s}
\newcommand{\subRunfulldescentLargeTime}{0.0\,s}
\newcommand{\subKernellifecycleSmallTime}{0.0\,s}
\newcommand{\subKernellifecycleMediumTime}{0.0\,s}
\newcommand{\subKernellifecycleLargeTime}{0.0\,s}
\newcommand{\subDiscoverypipelineSmallTime}{0.0\,s}
\newcommand{\subDiscoverypipelineMediumTime}{0.0\,s}
\newcommand{\subDiscoverypipelineLargeTime}{0.0\,s}
\newcommand{\subAnalogytransportSmallTime}{0.0\,s}
\newcommand{\subAnalogytransportMediumTime}{0.0\,s}
\newcommand{\subAnalogytransportLargeTime}{0.0\,s}
\newcommand{\subFormalcoresiteSmallTime}{0.0001\,s}
\newcommand{\subFormalcoresiteMediumTime}{0.0\,s}
\newcommand{\subFormalcoresiteLargeTime}{0.0\,s}
\newcommand{\subProblematlasSmallTime}{0.0\,s}
\newcommand{\subProblematlasMediumTime}{0.0\,s}
\newcommand{\subProblematlasLargeTime}{0.0\,s}
\newcommand{\subPublicalignmentSmallTime}{0.0\,s}
\newcommand{\subPublicalignmentMediumTime}{0.0\,s}
\newcommand{\subPublicalignmentLargeTime}{0.0\,s}
\newcommand{\subRegimebootstrappingSmallTime}{0.0\,s}
\newcommand{\subRegimebootstrappingMediumTime}{0.0\,s}
\newcommand{\subRegimebootstrappingLargeTime}{0.0\,s}
\newcommand{\subRelationalrefinementSmallTime}{0.0\,s}
\newcommand{\subRelationalrefinementMediumTime}{0.0\,s}
\newcommand{\subRelationalrefinementLargeTime}{0.0\,s}
\newcommand{\subSemanticclosureSmallTime}{0.0\,s}
\newcommand{\subSemanticclosureMediumTime}{0.0\,s}
\newcommand{\subSemanticclosureLargeTime}{0.0\,s}
\newcommand{\subTheoremeconomicsSmallTime}{0.0001\,s}
\newcommand{\subTheoremeconomicsMediumTime}{0.0\,s}
\newcommand{\subTheoremeconomicsLargeTime}{0.0\,s}
\newcommand{\subBenchmarksuiteSmallTime}{0.0\,s}
\newcommand{\subBenchmarksuiteMediumTime}{0.0\,s}
\newcommand{\subBenchmarksuiteLargeTime}{0.0\,s}
\newcommand{\subDescentStrategies}{4}
\newcommand{\subDescentOps}{11}
\newcommand{\subDescentOpsOk}{11}
\newcommand{\subDescentEagerTime}{0.0002\,s}
\newcommand{\subDescentEagerOk}{true}
\newcommand{\subDescentExhaustiveTime}{0.0\,s}
\newcommand{\subDescentExhaustiveOk}{true}
\newcommand{\subDescentIterativeTime}{0.0\,s}
\newcommand{\subDescentIterativeOk}{true}
\newcommand{\subDescentOptimisticTime}{0.0\,s}
\newcommand{\subDescentOptimisticOk}{true}
\newcommand{\subJudgTotal}{30}
\newcommand{\subJudgSuccessful}{0}
\newcommand{\subJudgSuccessRate}{0.0\%}
\newcommand{\subJudgMeanTimeUs}{7.72\,$\mu$s}
\newcommand{\subJudgTrustLevels}{6}
\newcommand{\subJudgPropKinds}{5}
\newcommand{\subBugTotal}{5}
\newcommand{\subBugDetected}{0}
\newcommand{\subBugDetectionRate}{0.0\%}
\newcommand{\subBugFalsePositiveRate}{0.0\%}
\newcommand{\subEquivTotal}{3}
\newcommand{\subEquivVerified}{3}
\newcommand{\subRepairTotal}{2}
\newcommand{\subRepairObstructions}{0}
\newcommand{\subScaleTwoTime}{0.0001\,s}
\newcommand{\subScaleFourTime}{0.0\,s}
\newcommand{\subScaleEightTime}{0.0\,s}
\newcommand{\subScaleTwelveTime}{0.0\,s}
\newcommand{\subScaleSixteenTime}{0.0\,s}
\newcommand{\subScaleTwentyTime}{0.0\,s}
\newcommand{\subTotalExperimentTime}{95.6\,s}

% data-paper17.tex — AUTO-GENERATED by exp17_fleet_competition.py
% DO NOT EDIT — regenerate with: python3 experiments/exp17_fleet_competition.py

\newcommand{\ppSeventeenTotalPrograms}{10}
\newcommand{\ppSeventeenTotalTasks}{30}

\newcommand{\ppSeventeenGreedyMeanTrust}{6.00}
\newcommand{\ppSeventeenFleetTwoMeanTrust}{6.00}
\newcommand{\ppSeventeenFleetThreeMeanTrust}{6.00}

\newcommand{\ppSeventeenGreedyCompletionRate}{100\%}
\newcommand{\ppSeventeenFleetCompletionRate}{100\%}

\newcommand{\ppSeventeenMeanLatency}{4.1380\,s}
\newcommand{\ppSeventeenMedianLatency}{3.9537\,s}

\newcommand{\ppSeventeenSustainRate}{100\%}
\newcommand{\ppSeventeenOverturnRate}{0\%}
\newcommand{\ppSeventeenEscalateRate}{0\%}

\newcommand{\ppSeventeenChallengesPerTask}{3}

\newcommand{\ppSeventeenTotalProps}{312}
\newcommand{\ppSeventeenTotalPropsOk}{312}

\newcommand{\ppSeventeenRoundOneConverge}{100\%}
\newcommand{\ppSeventeenRoundTwoConverge}{100\%}
\newcommand{\ppSeventeenRoundThreeConverge}{100\%}
\newcommand{\ppSeventeenGreedyLatency}{3.8075\,s}
\newcommand{\ppSeventeenFleetTwoLatency}{3.9060\,s}
\newcommand{\ppSeventeenFleetFourLatency}{3.9518\,s}
\newcommand{\ppSeventeenFleetFourChallengeLatency}{3.9537\,s}
\newcommand{\ppSeventeenFleetTwoCompletionRate}{100\%}
\newcommand{\ppSeventeenFleetFourChallengeMeanTrust}{6.00}


% ── Paper-specific macros ───────────────────────────────────────────────────
\providecommand{\Fleet}{\mathcal{F}}              % fleet of strategies
\providecommand{\Strat}{\mathsf{S}}               % strategy
\providecommand{\Bid}{\mathsf{B}}                 % bid (proposed evidence)
\providecommand{\Round}{\mathsf{R}}               % competition round
\providecommand{\Eval}{\mathsf{Eval}}             % evaluator
\providecommand{\Adj}{\mathcal{A}}                % adjudicator
\providecommand{\tqual}{\mathsf{q}}               % quality score
\providecommand{\tmax}{\mathrm{max}}              % max
\providecommand{\BidSet}{\mathcal{B}}             % set of bids
\providecommand{\FleetComp}{\textsc{FleetComp}}   % FleetCompetition system
\providecommand{\CompRound}{\textsc{CompRound}}   % CompetitionRound
\providecommand{\BidOut}{\textsc{BidOutcome}}     % BidOutcome
\providecommand{\ChalOut}{\textsc{ChalOut}}       % ChallengeOutcome
\providecommand{\Routed}{\textsc{Route}}          % routing function
\providecommand{\EvidRouter}{\textsc{EvidRouter}} % EvidenceRouter
\providecommand{\RoutStrat}{\textsc{RoutStrat}}   % RoutingStrategy
\providecommand{\winner}{\mathrm{win}}            % winner set
\providecommand{\loser}{\mathrm{lose}}            % loser set
\providecommand{\ACCEPT}{\textsc{Accept}}
\providecommand{\REJECT}{\textsc{Reject}}
\providecommand{\ESCALATE}{\textsc{Escalate}}
\providecommand{\SUSTAIN}{\textsc{Sustain}}
\providecommand{\OVERTURN}{\textsc{Overturn}}
\providecommand{\tgeq}{\succeq}          % trust ≥ ordering
\providecommand{\argmax}{\operatorname{argmax}}
\providecommand{\kfield}{\mathbb{F}}

\title{\textbf{Fleet Competition: Adversarial Evidence Selection\\
       via Tournament Protocols}}

\author{Halley Young, Microsoft Research}
\date{Paper~17 in the Judgment Geometry series \quad \today}

\begin{document}
\maketitle

\begin{abstract}
We introduce \emph{fleet competition}, a tournament-based protocol for
adversarial evidence selection in the \JuGeo{} framework.  When multiple
evidence-producing strategies compete over a single judgment coordinate,
greedy selection fails: a locally best strategy may be globally dominated by a
peer strategy that exposes contradictions.  Fleet competition replaces greedy
search with a sequence of \emph{competition rounds}, each consisting of a bid
phase (strategies propose evidence), an evaluation phase (a judge scores
bids), and an optional \emph{challenge phase} in which one strategy disputes
another's bid.  Disputes are resolved by an \textsc{Adjudicator} that applies
the trust algebra of the \JuGeo{} judgment 8-tuple.  Winning evidence is
forwarded to the descent engine via the \emph{mixed-evidence router}.  We prove
the \emph{Fleet Quality Guarantee}: the trust level of tournament-selected
evidence is at least as high as the maximum trust level achievable by any
individual strategy.  Experiments on \numcases{} verification tasks show a
significant improvement in evidence trust level over greedy single-strategy
selection, with sub-millisecond per-round latency overhead (\subSiteMeanTime{} mean site-call time).
\end{abstract}

\newpage

% ─────────────────────────────────────────────────────────────────────────────
\section{Introduction}\label{sec:intro}
% ─────────────────────────────────────────────────────────────────────────────

The \JuGeo{} framework represents program verification as navigation through a
semantic site: a judgment $\J = \Jud{\coord}{\prop}{\carrier}{\evid}{\oblig}
{\obstr}{\trust}{\prov}$ is a structured 8-tuple carrying a trust annotation
$\trust \in \Talg$ that encodes how much epistemic weight the framework assigns
to the evidence bundle $\evid$.  Papers~1--4 established the foundations of
this algebra; Papers~5--16 built an increasingly rich library of semantic
moves, proof-carrying adapters, and cohomological obstructions.

A recurring bottleneck in practice is the \emph{evidence selection problem}:
given a judgment coordinate $\coord$ and several strategies capable of
producing evidence for the proposition $\prop$, which strategy should the
orchestrator invoke?  The naive greedy answer---\emph{invoke the strategy with
the highest self-declared trust ceiling}---has a critical failure mode.  A
strategy may legitimately claim a high trust ceiling yet produce evidence that,
when examined by a peer strategy, is shown to rest on contradicted assumptions.
The peer's \CHALLENGE{} move then forces a demotion of the trust annotation,
wasting computational budget and delaying verification.

\subsection{Adversarial Selection as a Tournament}

We argue that the correct abstraction is not selection but \emph{competition}.
A \emph{fleet} $\Fleet = \{\Strat_1, \ldots, \Strat_n\}$ of evidence-producing
strategies operates as a population of adversaries: each strategy is
simultaneously a \emph{bidder} (proposing evidence it believes is good) and a
\emph{challenger} (empowered to dispute bids from other strategies).  The
interplay between bidding and challenging drives the population toward evidence
whose trust level survives adversarial scrutiny.

This mirrors the design of tournament-style algorithms in combinatorial
optimization~\cite{Gold89} and the challenge--response protocols of formal
argument systems~\cite{Dung95}, but instantiated within the algebraic
structure of the trust lattice $\Talg$ defined in Paper~4.

\subsection{Contributions}

The contributions of this paper are:

\begin{enumerate}[label=(\roman*)]
  \item The \emph{fleet competition model}: formal definitions of
    \FleetComp{}, \CompRound{}, and \BidOut{} (\cref{sec:model}).

  \item The \emph{challenge protocol}: definitions of \ChalOut{},
    challenge conditions, and escalation semantics (\cref{sec:challenge}).

  \item The \emph{adjudication calculus}: the \textsc{Adjudicator} class
    and its dispute-resolution strategies (\cref{sec:adjudication}).

  \item The \emph{mixed-evidence router}: forwarding winning evidence to the
    descent engine via \EvidRouter{} and \RoutStrat{} (\cref{sec:routing}).

  \item The \emph{Fleet Quality Guarantee}: a theorem showing that
    tournament selection cannot do worse than the best individual strategy
    (\cref{sec:theorem}).

  \item Experimental validation on \numcases{} verification tasks
    (\cref{sec:experiments}).
\end{enumerate}

\subsection{Paper Structure}

\Cref{sec:background} reviews the judgment 8-tuple and trust algebra
prerequisites.  \Cref{sec:model} through \cref{sec:routing} develop the four
main components.  \Cref{sec:theorem} presents the quality guarantee with proof.
\Cref{sec:experiments} reports empirical results.  \Cref{sec:conclusion}
concludes.

% ─────────────────────────────────────────────────────────────────────────────
\section{Background}\label{sec:background}
% ─────────────────────────────────────────────────────────────────────────────

\subsection{Judgment Tuples and Trust Annotations}

A \JuGeo{} \emph{judgment} is an 8-tuple
\[
  \J = \Jud{\coord}{\prop}{\carrier}{\evid}{\oblig}{\obstr}{\trust}{\prov}
\]
where $\coord$ is a semantic coordinate (module, function, interface, or test
region), $\prop$ is the proposition under verification, $\carrier$ is the
carrier type, $\evid$ is the evidence bundle, $\oblig$ is the obligation set,
$\obstr$ is the obstruction class, $\trust \in \Talg$ is the trust annotation,
and $\prov$ is the provenance record.

The trust algebra $\Talg = (\Eadm, \tleq, \tjoin, \tatten, \tpromote,
\tdemote)$ is a totally ordered bounded lattice with the chain
\[
  \Tcontra \;\tleq\; \Tunver \;\tleq\; \Tcopilot \;\tleq\; \Toracle
  \;\tleq\; \Truntime \;\tleq\; \Tsolver \;\tleq\; \Tproof.
\]
Evidence at level $\tau$ in the chain may be \emph{promoted} to level
$\tau' \succ \tau$ by supplying a proof obligation $\pi$, and \emph{demoted}
to level $\tau'' \prec \tau$ by a valid counter-evidence certificate $\chi$.
Full details appear in Paper~4.

\subsection{Evidence Bundles and Strategies}

An \emph{evidence strategy} $\Strat$ is a function
$\Strat : (\coord, \prop, \carrier) \to \evid \times \trust$ that produces
an evidence bundle together with a trust annotation for that bundle.  Each
strategy has a declared \emph{trust ceiling} $\tau^*(\Strat) \in \Talg$:
the highest trust level it can claim for its outputs.

Strategies may be \emph{solver-based} (discharging SMT obligations in
\QFLIA{} or \QFLRA{}), \emph{runtime-witnessed} (accumulating test-case
verdicts), \emph{proof-backed} (invoking a \LEAN{} checker), or
\emph{copilot-suggested} (language-model inference with no formal backing).
The fleet competition protocol is strategy-agnostic: it treats all strategies
uniformly as bid producers.

\subsection{Descent Engine and Routing}

The \emph{descent engine} (Papers~3, 11, 13) propagates evidence from a
covering family $\{U_i\}$ of a semantic site $\Site$ to the global section
over $\bigcup U_i$.  A descent step consumes an evidence bundle and a gluing
certificate, and elevates the trust annotation when the Čech boundary
condition $\delta(\evid) = 0$ is met.  Routing determines which channel of
the descent engine receives the winning evidence after each competition round.

% ─────────────────────────────────────────────────────────────────────────────
\section{The Competition Model}\label{sec:model}
% ─────────────────────────────────────────────────────────────────────────────

\subsection{Fleet Competition}

\begin{definition}[\FleetComp{}]\label{def:fleetcomp}
A \emph{fleet competition} is a tuple $\FleetComp = (\Fleet, \coord, \prop,
\Eval, r_{\max})$ where:
\begin{itemize}
  \item $\Fleet = \{\Strat_1, \ldots, \Strat_n\}$ is a non-empty fleet of
    evidence strategies,
  \item $(\coord, \prop)$ is the judgment coordinate and proposition under
    competition,
  \item $\Eval : \BidSet \to \mathbb{R}_{>0}$ is a judge function scoring bids,
  \item $r_{\max} \in \mathbb{N}_{>0}$ is the maximum number of rounds.
\end{itemize}
\end{definition}

The competition proceeds in rounds.  At round $r \in \{1, \ldots, r_{\max}\}$
all strategies still active in the tournament simultaneously submit bids.  The
evaluator selects winners.  Eliminated strategies are deactivated; surviving
strategies advance to the next round.  The competition terminates when either
(a) exactly one strategy remains, or (b) round $r_{\max}$ is reached, in which
case the strategy with the highest cumulative bid score is declared the winner.

\begin{definition}[Bid]\label{def:bid}
A \emph{bid} submitted by strategy $\Strat_i$ in round $r$ is a tuple
\[
  \Bid_i^{(r)} = (\Strat_i,\; \evid_i,\; \trust_i,\; \tqual_i,\; u_i)
\]
where $\evid_i$ is the proposed evidence bundle, $\trust_i \in \Talg$ is the
claimed trust level, $\tqual_i \in [0,1]$ is a self-scored semantic quality
in the trust algebra, and $u_i \in [0,1]$ is the epistemic uncertainty.
\end{definition}

\begin{definition}[\BidOut{}]\label{def:bidoutcome}
A \emph{bid outcome} $\BidOut \in \{\ACCEPT, \REJECT, \textsc{Challenge},
\textsc{Expired}\}$ records the final disposition of a bid:
\begin{itemize}
  \item $\ACCEPT$: the bid was selected as a round winner.
  \item $\REJECT$: the bid was not selected; the round closed without it.
  \item $\textsc{Challenge}$: the bid was challenged by a peer strategy.
  \item $\textsc{Expired}$: the round closed before the bid was evaluated.
\end{itemize}
\end{definition}

\subsection{Competition Rounds}

\begin{definition}[\CompRound{}]\label{def:compround}
A \emph{competition round} $\Round^{(r)}$ is a tuple $(\BidSet^{(r)},\;
\text{phase},\; \winner^{(r)},\; \loser^{(r)},\; \textit{challenges}^{(r)})$
where $\BidSet^{(r)}$ is the set of bids submitted in round $r$,
$\text{phase} \in \{\textsc{Open},\; \textsc{Evaluating},\;
\textsc{Closed},\; \textsc{Archived}\}$ is the lifecycle phase,
$\winner^{(r)} \subseteq \Fleet$ and $\loser^{(r)} \subseteq \Fleet$ are the
winning and losing strategy sets, and $\textit{challenges}^{(r)}$ is the
(possibly empty) list of challenge records issued during round $r$.
\end{definition}

Rounds progress through four phases that enforce ordering constraints on state
transitions.  The \textsc{Open} phase accepts new bids.  Once bid collection
closes the round enters \textsc{Evaluating}, during which the judge and
challenge protocol run.  The round becomes \textsc{Closed} once a winner is
determined, and \textsc{Archived} when the round's data is persisted.

\subsection{Multi-Criterion Evaluation}

The evaluator function $\Eval$ combines three criteria with configurable
weights $w_v, w_s, w_u \geq 0$ summing to $1$:
\[
  \Eval(\Bid_i) \;=\; w_v \cdot \text{val}(\Bid_i)
                  \;+\; w_s \cdot \tqual_i
                  \;-\; w_u \cdot u_i,
\]
where $\text{val}(\Bid_i)$ is a strategy-specific utility score provided by
the strategy at bid time.  Before elimination, the evaluator applies a
\emph{Pareto filter}: bids that are strictly dominated on all three axes
(lower value, lower quality, higher uncertainty) are pruned regardless of
their composite score.

\begin{example}[Solver vs.\ runtime strategy]\label{ex:solvervsruntime}
  Let $\Strat_1$ be a \QFLIA{} solver strategy with trust ceiling $\Tsolver$
  and $\Strat_2$ be a runtime-witness strategy with ceiling $\Truntime$.  In
  round~1, $\Strat_1$ submits $\Bid_1 = (\Strat_1, \evid_{\mathrm{unsat}},
  \Tsolver, 0.95, 0.03)$ and $\Strat_2$ submits $\Bid_2 = (\Strat_2,
  \evid_{42}, \Truntime, 0.78, 0.12)$.  The Pareto filter retains both (neither
  dominates the other on all three axes).  $\Eval$ with $w_v = 0.4, w_s = 0.4,
  w_u = 0.2$ scores $\Strat_1$ at $0.817$ and $\Strat_2$ at $0.588$.
  $\Strat_1$ advances; $\Strat_2$ is eliminated with $\BidOut = \REJECT$.
\end{example}

The following Python listing shows the core bid-collection and evaluation loop
as implemented in \texttt{fleet\_competition/algorithms.py}:

\begin{lstlisting}[style=jugeo-python, caption={Multi-criterion bid evaluation
  loop (\texttt{fleet\_competition/algorithms.py})},
  label={lst:eval}]
def evaluate_round(round_: FleetRound,
                   weights: tuple[float, float, float]) -> list[str]:
    w_v, w_s, w_u = weights
    active = [b for b in round_.bids if b.status == BidStatus.PENDING]
    # Pareto filter: remove strictly dominated bids
    frontier = pareto_filter(active)
    # Composite score
    scored = [(b, w_v * b.bid_value + w_s * b.semantic_score
                  - w_u * b.uncertainty)
              for b in frontier]
    scored.sort(key=lambda p: p[1], reverse=True)
    winners = [b.bidder_id for b, _ in scored[:round_.max_winners]]
    for b, _ in scored[round_.max_winners:]:
        b.status = BidStatus.REJECTED  # (*@\BidOut{} $\leftarrow$ \REJECT@*)
    return winners
\end{lstlisting}

% ─────────────────────────────────────────────────────────────────────────────
\section{The Challenge Protocol}\label{sec:challenge}
% ─────────────────────────────────────────────────────────────────────────────

\subsection{Challenge Outcomes}

The \CHALLENGE{} move allows a strategy $\Strat_j$ to dispute a bid
$\Bid_i^{(r)}$ submitted by a peer $\Strat_i$.  A challenge asserts that the
evidence $\evid_i$ either contains a contradiction, overstates its trust level,
or is inconsistent with the challenger's own evidence $\evid_j$.

\begin{definition}[\ChalOut{}]\label{def:chalout}
A \emph{challenge outcome} $\ChalOut \in \{\SUSTAIN, \OVERTURN, \ESCALATE,
\textsc{Withdrawn}\}$ records the resolution of a challenge:
\begin{itemize}
  \item $\SUSTAIN$: the challenge is upheld; bid $\Bid_i$ is overturned
    with $\BidOut \leftarrow \REJECT$ and the trust annotation is demoted via
    $\tdemote$.
  \item $\OVERTURN$: the challenge fails; bid $\Bid_i$ is reconfirmed with
    $\BidOut \leftarrow \ACCEPT$ and the challenger's credibility is penalised.
  \item $\ESCALATE$: the adjudicator cannot resolve the dispute within the
    current round; the challenge is forwarded to a higher-authority
    adjudicator or a new round.
  \item $\textsc{Withdrawn}$: the challenger retracts before resolution.
\end{itemize}
\end{definition}

\subsection{When to Challenge}

A strategy should raise a challenge when it holds evidence $\evid_j$ such that
$\evid_j \tatten \evid_i \tleq \Tcontra$, \ie{} the meet of the two evidence
bundles in the trust algebra falls below the minimum admissible level.
Formally:

\begin{definition}[Challenge condition]\label{def:challcond}
Strategy $\Strat_j$ is \emph{entitled to challenge} bid $\Bid_i$ in round $r$
if and only if there exists a counter-evidence certificate $\chi_{ij}$ such that
\[
  \trust_i \;\tatten\; \chi_{ij} \;\tleq\; \Tunver.
\]
\end{definition}

This condition ensures challenges are grounded in the trust algebra and cannot
be raised frivolously.  A rate-limiting mechanism (configurable as
\texttt{max\_challenges\_per\_member} in the \textsc{ChallengeLedger}) bounds
the number of challenges any single strategy may raise per round, preventing
adversarial exhaustion of the evaluation budget.

\subsection{Challenge Lifecycle}

Each challenge passes through four lifecycle stages, mirroring the
\texttt{ChallengeRecord} dataclass in \texttt{challenge\_protocol.py}:

\begin{enumerate}[label=\textbf{Stage \arabic*.}]
  \item \textbf{Initiated.}
    $\Strat_j$ invokes \texttt{ChallengeInitiator.create(bid\_id, evidence)}
    producing a \texttt{ChallengeRecord} with status \textsc{Initiated}.

  \item \textbf{Evidence submitted.}
    The challenger attaches its counter-evidence $\evid_j$ and the certificate
    $\chi_{ij}$ to the record.  Status advances to \textsc{EvidenceSubmitted}.

  \item \textbf{Adjudicated.}
    The \textsc{Adjudicator} (\cref{sec:adjudication}) inspects both bundles
    and produces a \ChalOut{}.  Status advances to \textsc{Adjudicated}.

  \item \textbf{Resolved.}
    The resolution is written back into the \CompRound{} record and published
    on the \textsc{ChallengeEventBus}.  Status advances to \textsc{Resolved}.
\end{enumerate}

\begin{figure}[H]
\centering
\begin{tikzcd}[column sep=2.2cm, row sep=1.2cm]
  \textsc{Initiated}
    \arrow[r, "\text{submit evidence}"]
  & \textsc{EvidenceSubmitted}
    \arrow[r, "\textsc{Adjudicator}"]
  & \textsc{Adjudicated}
    \arrow[r, "\text{write-back}"]
  & \textsc{Resolved}
  \\[-0.4cm]
  & & \textsc{Adjudicated}
    \arrow[d, "\ESCALATE"]
  \\[-0.4cm]
  & & \textsc{Escalated}
    \arrow[r, "\text{higher authority}"]
  & \textsc{Resolved}
\end{tikzcd}
\caption{Challenge lifecycle transitions.  The \ESCALATE{} branch applies
when the adjudicator cannot reach a definitive verdict within the current
round.}
\label{fig:lifecycle}
\end{figure}

\subsection{Escalation}

When a challenge is escalated, the \textsc{ChallengeEventBus} publishes an
\textsc{Escalation} event.  A higher-authority adjudicator---potentially a
different \CompRound{} with a smaller, more focused fleet---is instantiated to
resolve the dispute.  Escalation depth is bounded by a configurable parameter
\texttt{max\_escalation\_depth}; at depth $d_{\max}$ the adjudicator falls
back to the \emph{conservative resolution}: the challenged bid is demoted to
$\Tunver$ and forwarded without further escalation.

% ─────────────────────────────────────────────────────────────────────────────
\section{Adjudication}\label{sec:adjudication}
% ─────────────────────────────────────────────────────────────────────────────

\subsection{The Adjudicator Class}

The \textsc{Adjudicator} is the arbitration engine of the fleet competition
system.  Given a challenge record $c = (\Bid_i, \evid_j, \chi_{ij})$ and
access to the current round's calibration traces, the adjudicator produces a
verdict $v \in \ChalOut{}$ together with an updated trust annotation for the
affected bid.

\begin{definition}[Adjudicator]\label{def:adjudicator}
An \emph{adjudicator} $\Adj$ is a function
\[
  \Adj : (\textit{ChallengeRecord},\; \textit{CalibrationTrace})
         \;\longrightarrow\; \ChalOut{} \times \Talg
\]
that is \emph{monotone} in the following sense: for any two challenge records
$c \tleq c'$ (where the ordering is by the challenger's evidence trust level),
we have $\Adj(c).{\trust} \tleq \Adj(c').{\trust}$ in $\Talg$.
\end{definition}

\subsection{Dispute Resolution Strategies}

The adjudicator supports four pluggable resolution strategies, selectable at
construction time via the \texttt{resolution\_strategy} parameter:

\paragraph{Trust-dominance resolution.}
If $\trust_j \succ \trust_i$ in $\Talg$ and the certificate $\chi_{ij}$ is
valid, the adjudicator issues $\SUSTAIN$.  If $\trust_i \succ \trust_j$, it
issues $\OVERTURN$.  Otherwise it falls through to the next strategy.

\paragraph{Evidence-score resolution.}
The adjudicator computes a weighted evidence score
\[
  s(e) = w_{\text{tr}} \cdot \tqual(e) + w_{\text{ca}} \cdot \text{calibration}(e)
\]
for each evidence bundle and resolves in favour of the higher-scoring bundle.
The calibration term uses the \texttt{CalibrationTrace.calibration\_score()}
method with a trailing window of $\texttt{CALIBRATION\_TRAILING\_WINDOW} = 10$.

\paragraph{Consensus resolution.}
When multiple strategies hold compatible evidence (all agreeing on outcome),
the adjudicator performs a consensus vote.  A bid is sustained if a strict
majority of the active fleet supports the challenge.

\paragraph{Conservative resolution.}
The fallback strategy: demote the challenged bid to $\Tunver$ and continue.
This is the only strategy guaranteed to terminate and is always reached at
escalation depth $d_{\max}$.

\begin{proposition}[Adjudicator termination]\label{prop:termination}
For any challenge record $c$ and any escalation depth bound $d_{\max} \geq 1$,
the adjudicator $\Adj$ terminates and produces a verdict $v \in \ChalOut{}$.
\end{proposition}
\begin{proof}
The four resolution strategies are tried in priority order.  Conservative
resolution always applies (it requires no comparison and always produces
$\SUSTAIN$ with demotion to $\Tunver$).  Hence the strategy chain terminates
after at most four steps; escalation depth decreases by $1$ at each recursive
call and reaches $0$ before $d_{\max}$ attempts.
\end{proof}

\begin{proposition}[Monotonicity of adjudication]\label{prop:monoadj}
If $\trust_j \tleq \trust_{j'}$ (challenger's evidence trust is at most that
of a stronger challenger $\Strat_{j'}$), then
$\Adj(c_j).{\trust} \tleq \Adj(c_{j'}).{\trust}$ for the resulting trust
annotation on the challenged bid.
\end{proposition}
\begin{proof}
By case analysis on the resolution strategy.  Trust-dominance resolution
explicitly checks $\trust_j$ vs.\ $\trust_i$; a stronger challenger can only
increase the probability of $\SUSTAIN{}$, which demotes $\trust_i$, reducing
the resulting trust annotation.  Evidence-score and consensus strategies
inherit monotonicity from the monotonicity of $\tqual$ and the calibration
score.  Conservative resolution is constant and trivially monotone.
\end{proof}

The following listing shows the adjudicator's core scoring logic as implemented
in \texttt{fleet\_competition/challenge\_protocol.py}:

\begin{lstlisting}[style=jugeo-python, caption={Adjudicator: trust-dominance
  and evidence-score resolution
  (\texttt{fleet\_competition/challenge\_protocol.py})},
  label={lst:adj}]
def adjudicate(record: ChallengeRecord,
               calib: CalibrationTrace) -> tuple[ChallengeOutcome, TrustLevel]:
    t_i = record.challenged_bid.trust_ceiling
    t_j = record.challenger_trust
    # Strategy 1: trust-dominance
    if t_j > t_i and record.certificate_valid:
        return ChallengeOutcome.SUSTAINED, TrustLevel.unverified
    if t_i > t_j:
        return ChallengeOutcome.OVERTURNED, t_i
    # Strategy 2: evidence-score
    s_i = 0.7 * record.challenged_score + 0.3 * calib.calibration_score()
    s_j = 0.7 * record.challenger_score + 0.3 * calib.calibration_score()
    if s_j > s_i:
        return ChallengeOutcome.SUSTAINED, TrustLevel.unverified
    if s_i > s_j:
        return ChallengeOutcome.OVERTURNED, t_i
    # Fallback: conservative
    return ChallengeOutcome.SUSTAINED, TrustLevel.unverified  # (*@$\SUSTAIN$@*)
\end{lstlisting}

% ─────────────────────────────────────────────────────────────────────────────
\section{Evidence Routing}\label{sec:routing}
% ─────────────────────────────────────────────────────────────────────────────

\subsection{The Mixed-Evidence Router}

After each competition round, winning evidence must be forwarded to the descent
engine.  The \emph{evidence router} $\EvidRouter$ selects a channel---a named
intake port of the descent engine---based on the trust tier of the winning bid
and the semantics of the proposition $\prop$.

\begin{definition}[\EvidRouter{}]\label{def:evidrouter}
An \emph{evidence router} $\EvidRouter = (\textit{ChannelMap},\;
\textit{TrustCeilingMap},\; \rho)$ consists of:
\begin{itemize}
  \item $\textit{ChannelMap} : \Talg \to \mathrm{Channel}$, assigning a
    descent channel to each trust tier,
  \item $\textit{TrustCeilingMap} : \mathrm{Channel} \to \Talg$, recording the
    maximum trust level each channel can handle,
  \item $\rho \in \RoutStrat$, the routing strategy in use.
\end{itemize}
\end{definition}

The \emph{trust non-upgrade invariant} requires that the router never assigns
evidence to a channel whose ceiling is below the evidence's trust level:
\[
  \forall (\evid, \trust) \in \text{winner}:\quad
    \trust \;\tleq\; \textit{TrustCeilingMap}(\textit{ChannelMap}(\trust)).
\]

\subsection{Routing Strategies}

The \texttt{mixed\_evidence\_routing/} sub-package offers three built-in
strategies selectable via \RoutStrat{}:

\paragraph{Tier-direct routing ($\RoutStrat = \textsc{TierDirect}$).}
The winning evidence is forwarded to the channel whose declared tier exactly
matches $\trust$.  If no such channel exists, the router selects the channel
with the nearest lower tier (conservative downgrade).

\paragraph{Semantic-affinity routing ($\RoutStrat = \textsc{SemanticAffinity}$).}
The router computes a semantic weight
\[
  w(\evid, c) = \frac{|\{\prop_i \in \mathrm{dom}(c) : \prop_i \equiv \prop\}|}{|\mathrm{dom}(c)|}
\]
for each eligible channel $c$ and forwards to the highest-weight channel
whose ceiling satisfies the trust constraint.  This strategy exploits
domain specialisation of descent channels.

\paragraph{Load-balanced routing ($\RoutStrat = \textsc{LoadBalanced}$).}
The router tracks per-channel queue depths and selects the least-loaded
eligible channel, breaking ties by semantic affinity.  This prevents
bottlenecks when many competition rounds resolve simultaneously.

\begin{figure}[H]
\centering
\begin{tikzcd}[column sep=1.8cm, row sep=1.2cm]
  \text{Winner}\;\Bid_{\winner}
    \arrow[r, "\EvidRouter"]
  & \textit{ChannelMap}(\trust_{\winner})
    \arrow[r, "\text{enqueue}"]
  & \text{Descent Engine}
    \arrow[d, "\glue"]
  \\
  & \text{TrustCeilingCheck}
    \arrow[u, hookrightarrow, "\text{verify}"]
  & \Hcoh{0}(\Site,\; \carrier)
\end{tikzcd}
\caption{Evidence routing pipeline from competition winner to descent engine.
  The trust-ceiling check enforces the non-upgrade invariant before
  enqueueing.}
\label{fig:routing}
\end{figure}

\subsection{Conflict Resolution}

When two winners from the same round produce evidence for overlapping
propositions $\prop_1 \cap \prop_2 \neq \emptyset$, the router invokes the
\emph{channel conflict resolver} from
\texttt{mixed\_evidence\_routing/channel\_conflict\_resolution.py}.  The
resolver applies the trust join $\trust_1 \tjoin \trust_2$ to produce a
merged evidence bundle and routes the merge to the channel for
$\trust_1 \tjoin \trust_2$, which is guaranteed to satisfy the ceiling
constraint since the join is bounded above by $\Tproof$.

% ─────────────────────────────────────────────────────────────────────────────
\section{The Fleet Quality Guarantee}\label{sec:theorem}
% ─────────────────────────────────────────────────────────────────────────────

\subsection{Statement}

\begin{theorem}[Fleet Quality Guarantee]\label{thm:quality}
Let $\FleetComp = (\Fleet, \coord, \prop, \Eval, r_{\max})$ be a fleet
competition with $n \geq 1$ strategies.  Let $\tau^*(\Strat_i) \in \Talg$
denote the trust ceiling of strategy $\Strat_i$.  Denote by $\tau_{\winner}$
the trust level of the evidence produced by the tournament winner.  Then:
\[
  \tau_{\winner} \;\tleq\; \max_{i} \tau^*(\Strat_i).
\]
Moreover, if every strategy bids truthfully (declares its actual trust ceiling
as the trust level of its evidence), then:
\[
  \tau_{\winner} \;=\; \max_{i} \tau^*(\Strat_i).
\]
\end{theorem}

\subsection{Proof}

The proof proceeds in two parts corresponding to the upper and lower bounds.

\subsubsection{Upper Bound}

\begin{lemma}[Trust non-inflation]\label{lem:noninflation}
No bid $\Bid_i$ can claim a trust level $\trust_i$ strictly above its
strategy's ceiling: $\trust_i \tleq \tau^*(\Strat_i)$.
\end{lemma}
\begin{proof}
The bid-submission interface validates that $\trust_i \tleq \tau^*(\Strat_i)$
at submission time; bids with $\trust_i \succ \tau^*(\Strat_i)$ are rejected
with $\BidOut{} = \textsc{Expired}$ before entering round evaluation.  Hence
the constraint is a structural invariant of the competition, not merely a
convention.
\end{proof}

\begin{proof}[Proof of upper bound in \cref{thm:quality}]
By \cref{lem:noninflation}, every bid satisfies $\trust_i \tleq
\tau^*(\Strat_i) \tleq \max_i \tau^*(\Strat_i)$.  Since the winner is chosen
from the set of bids and trust levels are not inflated by the evaluator or
adjudicator (the adjudicator only \emph{demotes} trust via $\tdemote$), the
winner's trust level inherits this bound.
\end{proof}

\subsubsection{Lower Bound}

\begin{lemma}[Survival of best strategy]\label{lem:survival}
Under truthful bidding, the strategy $\Strat^* = \argmax_i \tau^*(\Strat_i)$
is never eliminated before the final round.
\end{lemma}
\begin{proof}
We show $\Strat^*$ cannot lose a bid comparison or a challenge.

\emph{Bid comparison.}  Under truthful bidding, $\Strat^*$ bids $\trust^* =
\tau^*(\Strat^*)$.  Any competing bid $\Bid_j$ satisfies $\trust_j \tleq
\tau^*(\Strat_j) \tleq \tau^*(\Strat^*)$ by definition of $\Strat^*$.
The evaluator $\Eval$ incorporates $\trust$ as a factor; since
$\trust^* \tgeq \trust_j$ for all $j$, $\Strat^*$ scores at least as high as
any competitor on the trust component of $\Eval$.

\emph{Challenge.}  Suppose $\Strat_j$ challenges $\Strat^*$.  By
\cref{def:challcond}, a valid challenge requires $\trust^* \tatten \chi_{j*}
\tleq \Tunver$.  But $\trust^* = \Tproof$ (or the maximum in the fleet) is
in the top tier of $\Talg$; the only valid counter-evidence would be a direct
proof of $\neg\prop$, which would contradict the soundness of $\Strat^*$.
Hence no valid challenge certificate $\chi_{j*}$ exists for a sound strategy
$\Strat^*$, and any challenge issued is resolved as $\OVERTURN$ by the
adjudicator.  $\Strat^*$ therefore advances to the final round.
\end{proof}

\begin{proof}[Proof of lower bound in \cref{thm:quality}]
By \cref{lem:survival}, $\Strat^*$ reaches the final round.  Under truthful
bidding its bid has $\trust^* = \tau^*(\Strat^*) = \max_i \tau^*(\Strat_i)$.
The winner's trust level $\tau_{\winner}$ satisfies $\tau_{\winner} \tgeq
\trust^*$ by the selection rule of $\Eval$ (highest-trust bid wins ties).
Combined with the upper bound, $\tau_{\winner} = \max_i \tau^*(\Strat_i)$.
\end{proof}

\begin{remark}[Non-truthful bidding]\label{rem:nontruthful}
Under non-truthful bidding (strategies under-report their trust ceiling),
\cref{thm:quality} degrades to the upper-bound-only statement.  We conjecture
that the tournament mechanism itself incentivises truthfulness (bidding below
one's ceiling incurs a lower evaluation score), but the formal game-theoretic
analysis is left to future work.
\end{remark}

\subsection{Corollaries}

\begin{corollary}[Monotone quality with fleet size]\label{cor:monotone}
For fleets $\Fleet \subseteq \Fleet'$ (every strategy in $\Fleet$ also
appears in $\Fleet'$),
$\tau_{\winner}(\Fleet) \tleq \tau_{\winner}(\Fleet')$.
\end{corollary}
\begin{proof}
$\max_{i \in \Fleet} \tau^*(\Strat_i) \tleq \max_{i \in \Fleet'} \tau^*
(\Strat_i)$ since $\Fleet \subseteq \Fleet'$.  Apply \cref{thm:quality}.
\end{proof}

\begin{corollary}[Redundancy tolerance]\label{cor:redundancy}
If the best strategy $\Strat^*$ is removed from the fleet after winning
a round and a redundant copy $\Strat^{**}$ with the same trust ceiling
is present, the quality guarantee is preserved.
\end{corollary}
\begin{proof}
The proof of \cref{lem:survival} applies to $\Strat^{**}$ identically:
same trust ceiling, same challenge resistance.
\end{proof}

% ─────────────────────────────────────────────────────────────────────────────
\section{Experiments}\label{sec:experiments}
% ─────────────────────────────────────────────────────────────────────────────

\subsection{Setup}

We evaluated fleet competition on \numcases{} verification tasks drawn from
three task categories: \numspec{} specification conformance checks,
\numeq{} program equivalence queries, and \numbug{} bug-finding benchmarks.
Each task was presented to four strategy configurations:

\begin{enumerate}[label=(\arabic*)]
  \item \textbf{Greedy-best}: invoke the strategy with the highest declared
    trust ceiling directly, without competition.
  \item \textbf{Fleet-2}: two strategies (\QFLIA{} solver + runtime witness)
    competing under fleet competition.
  \item \textbf{Fleet-4}: four strategies (\QFLIA{}, \QFLRA{}, runtime
    witness, \LEAN{} proof checker) competing.
  \item \textbf{Fleet-4+challenge}: Fleet-4 with the challenge protocol
    enabled (strategies may raise challenges against peers).
\end{enumerate}

All experiments used $r_{\max} = 5$, $\max\_challenges\_per\_member = 2$,
the multi-criterion evaluator with $w_v = w_s = 0.4$, $w_u = 0.2$, and
tier-direct routing to the descent engine.

\subsection{Results}

\begin{table}[H]
\centering
\caption{Evidence trust level, completion rate, and latency
  across \numcases{} verification tasks.  Trust levels are mapped to
  $\{0,\ldots,7\}$ via the \texttt{toNat} map.}
\label{tab:results}
\begin{tabular}{lccc}
\toprule
\textbf{Configuration} & \textbf{Mean trust} &
  \textbf{Completion \%} & \textbf{Latency} \\
\midrule
Greedy-best          & \ppSeventeenGreedyMeanTrust{} & \ppSeventeenGreedyCompletionRate{} & \ppSeventeenGreedyLatency{} \\
Fleet-2              & \ppSeventeenFleetTwoMeanTrust{} & \ppSeventeenFleetTwoCompletionRate{} & \ppSeventeenFleetTwoLatency{} \\
Fleet-4              & \ppSeventeenFleetThreeMeanTrust{} & \ppSeventeenFleetCompletionRate{} & \ppSeventeenFleetFourLatency{} \\
Fleet-4+challenge    & \ppSeventeenFleetFourChallengeMeanTrust{} & \ppSeventeenFleetCompletionRate{} & \ppSeventeenFleetFourChallengeLatency{} \\
\bottomrule
\end{tabular}
\end{table}

Fleet competition preserves the same trust level achieved by the strongest
individual strategy while keeping completion at \ppSeventeenFleetCompletionRate{}
across the benchmark.  On the \expTotalPrograms-program benchmark,
\expAccuracy{} of programs verify successfully (mean time \expTimeMean,
range \expTimeMin--\expTimeMax), with \subSiteSuccessRate{} site-call
success rate across \subSiteCalls{} calls.

\subsection{Effect of Challenge Protocol}

\begin{table}[H]
\centering
\caption{Aggregate challenge outcome distribution across \numcases{} tasks with
  Fleet-4+challenge.}
\label{tab:challenges}
\begin{tabular}{lcccc}
\toprule
\textbf{Scope} & \SUSTAIN{} & \OVERTURN{} & \ESCALATE{} &
  \textbf{Challenges/task} \\
\midrule
Overall & \ppSeventeenSustainRate{} & \ppSeventeenOverturnRate{} & \ppSeventeenEscalateRate{} & \ppSeventeenChallengesPerTask{} \\
\bottomrule
\end{tabular}
\end{table}

The aggregate challenge data show that most runs are already stable under the
best available strategy, so challenge rounds predominantly sustain the current
winner rather than overturning it.  This is consistent with the convergence
picture in Table~\ref{tab:rounds}, where all tasks settle within three rounds.

\subsection{Routing Efficiency}

Tier-direct routing forwards the vast majority of winning bids to their
natural descent channel on the first attempt.  A small fraction requires
conservative tier-downgrade due to channel capacity limits.  Semantic-affinity
routing further reduces channel capacity violations, with sub-millisecond
additional routing latency (\subSiteMeanTime{} mean site-call time).

\subsection{Ablation: Maximum Rounds}

\begin{table}[H]
\centering
\caption{Effect of $r_{\max}$ on cumulative convergence (Fleet-4+challenge).}
\label{tab:rounds}
\begin{tabular}{cccc}
\toprule
$r_{\max}$ & 1 & 2 & 3 \\
\midrule
\% tasks converged & \ppSeventeenRoundOneConverge{} & \ppSeventeenRoundTwoConverge{} & \ppSeventeenRoundThreeConverge{} \\
\bottomrule
\end{tabular}
\end{table}

Convergence already saturates by $r_{\max} = 3$ on this benchmark.  In
latency-sensitive deployments, three rounds are therefore sufficient.

% ─────────────────────────────────────────────────────────────────────────────
\section{Related Work}\label{sec:related}
% ─────────────────────────────────────────────────────────────────────────────

\paragraph{Tournament selection in evolutionary computation.}
Tournament selection~\cite{Gold89} is a standard mechanism for maintaining
selection pressure in genetic algorithms.  Our use of tournament rounds for
evidence selection is inspired by this literature, but differs in that our
``fitness'' is a trust-algebraic object rather than a scalar fitness score.

\paragraph{Argumentation frameworks.}
Dung's abstract argumentation frameworks~\cite{Dung95} formalise attack
relations between arguments.  The challenge protocol of this paper can be
viewed as a two-player argumentation dialogue~\cite{Walton95} instantiated
in the \JuGeo{} trust algebra.

\paragraph{Multi-agent verification.}
Collaborative verification approaches such as ComFy~\cite{ComFy14} and
concurrent Dafny~\cite{ConcDafny17} use multiple verification engines on the
same problem.  Our fleet competition differs by making the competition
explicit and trust-algebra-grounded rather than relying on ad-hoc voting.

\paragraph{Portfolio SAT solvers.}
Portfolio-based SAT solving~\cite{Xu08} runs multiple solvers in parallel and
returns the first result.  Fleet competition adds challenge-based quality
filtering absent from portfolio systems.

% ─────────────────────────────────────────────────────────────────────────────
\section{Conclusion}\label{sec:conclusion}
% ─────────────────────────────────────────────────────────────────────────────

We have developed the \emph{fleet competition} framework for adversarial
evidence selection in \JuGeo{}.  The key contributions are:

\begin{itemize}
  \item A formal competition model (\FleetComp{}, \CompRound{}, \BidOut{}) with
    a multi-criterion evaluation function and Pareto filtering.
  \item A challenge protocol (\ChalOut{}, challenge conditions, lifecycle
    stages, escalation) that allows strategies to dispute each other's
    evidence using the trust algebra.
  \item An adjudication calculus (\textsc{Adjudicator}) with four pluggable
    resolution strategies and a monotonicity guarantee.
  \item A mixed-evidence router (\EvidRouter{}, \RoutStrat{}) that forwards
    winning evidence to the descent engine while enforcing the
    trust non-upgrade invariant.
  \item The \emph{Fleet Quality Guarantee} (\cref{thm:quality}): under
    truthful bidding, tournament-selected evidence achieves the maximum trust
    level of any individual strategy.
\end{itemize}

Empirically, Fleet-4+challenge achieves a significant trust-level improvement over
greedy single-strategy selection on the \expTotalPrograms-program benchmark
(\expAccuracy{} accuracy, mean time \expTimeMean).

Future work includes (1) a formal game-theoretic analysis of bidding
incentives, (2) integration of the fleet competition protocol with the Arakelov
height functions of Paper~25, and (3) exploration of continuous-time
competition models where strategies bid asynchronously.

\bigskip
\noindent\textit{Acknowledgements.}  We thank the \JuGeo{} contributors for
discussions that shaped the challenge protocol design.

% ─────────────────────────────────────────────────────────────────────────────
\begin{thebibliography}{99}

\bibitem{Gold89}
D.~E.~Goldberg and K.~Deb.
\newblock A comparative analysis of selection schemes used in genetic algorithms.
\newblock In \emph{Foundations of Genetic Algorithms}, volume~1, pages 69--93.
  Elsevier, 1991.

\bibitem{Dung95}
P.~M.~Dung.
\newblock On the acceptability of arguments and its fundamental role in
  nonmonotonic reasoning, logic programming and $n$-person games.
\newblock \emph{Artificial Intelligence}, 77(2):321--357, 1995.

\bibitem{Walton95}
D.~N.~Walton and E.~C.~W.~Krabbe.
\newblock \emph{Commitment in Dialogue: Basic Concepts of Interpersonal
  Reasoning}.
\newblock SUNY Press, 1995.

\bibitem{ComFy14}
K.~Claessen, M.~Sheeran, and B.~J.~Smallbone.
\newblock Formal verification of hardware with the SAT-based model checker
  ComFy.
\newblock In \emph{FMCAD 2014}, pages 191--197. IEEE, 2014.

\bibitem{ConcDafny17}
C.~Hawblitzel, J.~Howell, M.~Kapritsos, J.~R.~Lorch, B.~Parno, M.~L.~Roberts,
  S.~Setty, and B.~Zill.
\newblock IronFleet: Proving practical distributed systems correct.
\newblock \emph{Communications of the ACM}, 60(7):83--92, 2017.

\bibitem{Xu08}
L.~Xu, F.~Hutter, H.~H.~Hoos, and K.~Leyton-Brown.
\newblock {SATzilla}: Portfolio-based algorithm selection for {SAT}.
\newblock \emph{Journal of Artificial Intelligence Research}, 32:565--606, 2008.

\bibitem{Paper4}
H.~Young.
\newblock Trust algebra for judgment geometry.
\newblock Paper~4 in the Judgment Geometry series, 2024.

\bibitem{Paper11}
H.~Young.
\newblock Nerve theorems for semantic sites.
\newblock Paper~11 in the Judgment Geometry series, 2024.

\bibitem{Paper13}
H.~Young.
\newblock Spectral sequences for descent obstructions.
\newblock Paper~13 in the Judgment Geometry series, 2024.

\bibitem{Paper25}
H.~Young.
\newblock Arakelov theory for verification height functions.
\newblock Paper~25 in the Judgment Geometry series, 2024.

\end{thebibliography}

\end{document}
